\subsection{Определение этапов работ}

Для разработки программного средства обнаружения аномалий в Kubernetes-кластере на основе анализа audit-логов необходимо определить последовательные этапы выполнения работ, позволяющие структурировать процесс технической реализации проекта. В качестве основы для выделения стадий разработки используется ГОСТ 19.102-77 «Стадии разработки» \cite{gost_19_102_77}, который регламентирует порядок создания программ и программной документации независимо от области их применения.

Стадия \textit{«Техническое задание»} включает следующие этапы:
\begin{enumerate}
    \item обоснование необходимости разработки ПО;
    \item проведение научно-исследовательской работы;
    \item разработка и утверждение технического задания.
\end{enumerate}

Данная стадия необходима для постановки целей и задач проекта, анализа предметной области, изучения возможных подходов к обработке audit-логов, определения требований к функциональности и качеству программного средства, а также формирования технического задания на разработку.

Стадия \textit{«Эскизный проект»} состоит из следующих этапов:
\begin{enumerate}
    \item разработка эскизного проекта;
    \item утверждение эскизного проекта.
\end{enumerate}

На данном этапе уточняются методы решения поставленной задачи, формируется общее описание алгоритмов анализа событий в Kubernetes-кластере, определяется структура входных и выходных данных, а также подготавливается пояснительная записка с описанием проектных решений.

Стадия \textit{«Технический проект»} состоит из следующих этапов:
\begin{enumerate}
    \item разработка технического проекта;
    \item утверждение технического проекта.
\end{enumerate}

На этом этапе выполняется разработка алгоритмов обработки audit-логов, проектирование структуры программного средства, определение модулей системы, а также планирование последующей реализации и тестирования комплекса.

Стадия \textit{«Рабочий проект»} содержит следующие этапы:
\begin{enumerate}
    \item разработка программы;
    \item разработка программной документации;
    \item испытание программы.
\end{enumerate}

На данной стадии осуществляется непосредственная реализация программного комплекса, его отладка, тестирование и корректировка. Параллельно разрабатывается комплект программной документации, включающий описание работы системы, инструкции по развертыванию и сопровождению. В соответствии с результатами прохождения испытаний и приёмо-сдаточных мероприятий производится корректировка работы программы и внесение изменений в документацию.

Стадия \textit{«Внедрение»} заключается в подготовке программного средства к передаче заказчику, его установке в целевой инфраструктуре и организации дальнейшего сопровождения.

Таким образом, в рамках разработки программного средства обнаружения аномалий в Kubernetes-кластере на основе анализа audit-логов можно выделить следующие основные этапы:

\begin{enumerate}
    \item исследовательская часть;
    \item разработка архитектуры системы;
    \item разработка и исследование алгоритмов НС;
    \item анализ и выбор наилучшей модели НС;
    \item разработка программных модулей;
    \item интеграция выбранной модели в программный комплекс;
    \item тестирование и отладка программного средства;
    \item разработка технической документации.
\end{enumerate}
